% !TEX root = ../main.tex

\chapter{实验系统}

本章介绍横向射流火焰实验平台、双喷嘴构型及声激励系统，以及温度测量系统。

\section{横向射流火焰实验平台}

本研究搭建的横向射流火焰实验台采用\textbf{两级燃烧室串联}设计，其原理示意图如图~\ref{fig:10}所示，搭建实验台的实物图如图~\ref{fig:11}所示。在该平台中，左侧通入的主流燃料与空气在第一级燃烧室内形成稳定的多重扩散火焰。其高温产物气流经过一个收缩段实现流场均匀化后，进入二级燃烧室。在此处通入次级燃料射流，高温主流即将其点燃，形成横向射流火焰。实验首先明确了包括燃料类型、横向射流速度及喷射压力在内的基础工况参数及其可调范围，并建立了基于高速摄影与温度测量的火焰宏观结构、温度分布及燃烧效率的系统化分析方法。

\textbf{一级燃烧室}：通入主流燃料（甲烷）与空气，形成稳定的多重扩散火焰，产生高温（1200-1500~K）产物气流。

\textbf{收缩段}：将高温气流均匀化，确保进入二级燃烧室的来流速度分布均匀。

\textbf{二级燃烧室}：通过底部的双喷嘴注入次级燃料与空气，高温主流将其点燃，形成横向射流火焰。燃烧室采用\textbf{石英玻璃材质}，便于光学诊断与高速摄影观测。

实验台整体采用欧标40×40铝型材框架支撑，便于移动实验台的位置，方便进行实验图片的拍摄与光学对准。

\begin{figure}[!htbp]
	\centering
	\includegraphics[width=0.9\textwidth]{image10.png}
	\caption{横向射流火焰实验台原理示意图}
	\label{fig:10}
\end{figure}

\begin{figure}[!htbp]
	\centering
	\includegraphics[width=0.9\textwidth]{image11.jpg}
	\caption{横向射流火焰实验台实物图}
	\label{fig:11}
\end{figure}

\section{双喷嘴设计与声激励系统}

\subsection{双喷嘴构型设计}

为实现燃料与空气的独立调控与精准混合，本研究采用\textbf{双喷嘴构型}：前喷嘴通入燃料（甲烷），后喷嘴通入空气。两种气体在到达二级燃烧室出口后才开始混合，避免预混带来的安全隐患。截至目前，已完成喷嘴加工工艺优化、多规格双喷嘴图纸绘制，均已加工完成得到实物。

喷嘴加工要求提出与质量检验：针对双喷嘴的加工，采用开孔插入式焊接工艺，即开出与细管外径大小相同的孔，将细管插入后再对其进行焊接，为保证输入的气体效果不受加工工艺的影响，要求焊接处维持圆管形状、内部打磨光滑。通过多次质量检验和反馈，最终得到可以使用的双喷嘴加工件用于实验当中。

双喷嘴尺寸评估与加工方案调整：根据导师意见，确认现有6~mm内径喷嘴偏大。确立了新图纸加工小喷嘴+现有喷嘴临时使用的方案，即利用现有6~mm喷嘴进行预实验和系统调试，并对工况进行相应的换算，同时加工2~mm、3~mm内径的小喷嘴用于正式实验。绘制了4组不同内径和间距的双喷嘴图纸（2~mm内径4~mm间距、2~mm内径6~mm间距、3~mm内径6~mm间距、3~mm内径9~mm间距），明确了316L不锈钢材质、外部抛光、内壁光滑、密封面粗糙度Ra≤1.6等详细加工要求，并得到可用的双喷嘴加工件。双喷嘴的参数如表~\ref{tab:2}所示。

\begin{table}[!htbp]
	\centering
	\caption{双喷嘴参数表}
	\label{tab:2}
	\begin{tabular}{c|c|c}
		\hline
		编号 & 孔径$d$（mm） & 间距$s$ \\
		\hline
		a & 2 & 2d（4mm） \\
		b & 2 & 3d（6mm） \\
		c & 3 & 2d（6mm） \\
		d & 3 & 3d（9mm） \\
		\hline
	\end{tabular}
\end{table}

图~\ref{fig:12}、图~\ref{fig:13}展示了2~mm和3~mm孔径双喷嘴的工程图纸，图~\ref{fig:14}为加工完成的实物照片。加工采用\textbf{开孔插入式焊接工艺}，确保焊接处维持圆管形状、内部平滑。

\begin{figure}[!htbp]
	\centering
	\includegraphics[width=0.8\textwidth]{image12.jpg}
	\caption{2~mm双喷嘴图纸}
	\label{fig:12}
\end{figure}

\begin{figure}[!htbp]
	\centering
	\includegraphics[width=0.8\textwidth]{image13.jpg}
	\caption{3~mm双喷嘴图纸}
	\label{fig:13}
\end{figure}

\begin{figure}[!htbp]
	\centering
	\includegraphics[width=0.4\textwidth]{image14.jpg}
	\caption{双喷嘴实物图}
	\label{fig:14}
\end{figure}

\subsection{声激励脉动系统}

为精确引入可控的上游速度脉动，本研究设计并集成了\textbf{声激励脉动系统}。系统以\textbf{扬声器}为核心声源，通过刚性腔体与细管将声能定向传导至喷嘴上游区域。

\textbf{系统关键参数：}
\begin{itemize}
	\item 频率范围：45~Hz--2.5~kHz
	\item 波形：正弦波、方波等（其中方波被证实有更强的离散涡结构）
	\item 幅值：可通过信号发生器进行调节
	\item 气路设计：甲烷与空气分别独立设置声激励罐体，扬声器安装于罐体尾部，气体从侧壁引入，避免气流对扬声器膜片的冲击
\end{itemize}

声激励喷嘴初步设计：根据声激励脉动系统的所需功能对其总体构造进行设计，确定其由尾部法兰、内部法兰、尾部密封不锈钢板等主要部件构成，其整体结构示意图如图~\ref{fig:15}所示。关键创新点在于使甲烷与空气分别进入二级燃烧室当中，使燃料在双喷嘴出口处与空气混合，避免预混导致危险情况的发生。同时确定了扬声器、1000~W电源、信号发生板等关键元器件的参数，形成了完整的采购清单。

\begin{figure}[!htbp]
	\centering
	\includegraphics[width=0.5\textwidth]{image15.png}
	\caption{声激励脉动系统图纸}
	\label{fig:15}
\end{figure}

扬声器与电路初步测试：完成了扬声器的发声测试，验证了其在罐体内的工作性能。测试中发现了密封不严、电源隔离不足、布线混乱等问题，针对性地提出了采购4~mm厚垫片、单独配置保护电源、改用双芯线和三芯线等解决方案，并最终解决其出现的问题。

扬声器安装与腔体改造：与课题组人员协作完成了扬声器、航插、气路的安装固定。通过截短喷嘴细管、焊接空气通入管，使罐体满足气路和电路连接要求。针对双喷嘴实验需求，预留了电路空间和管路接口。并对电路重新设计与布局优化，采用接线端子替换原有散线，重新规划了电路连接方案，显著改善了电路布局的整洁性和安全性，降低了电路故障风险。

罐体总装与密封测试：完成了两套罐体的整体组装，分别用于甲烷和空气气路，确定了上游设置流量计、垫片密封底部的方案，利用流量计的限压作用避免罐体承压过大导致密封破坏。实物图如图~\ref{fig:16}所示。

\begin{figure}[!htbp]
	\centering
	\includegraphics[width=0.5\textwidth]{image16.jpg}
	\caption{声激励脉动系统实物图}
	\label{fig:16}
\end{figure}

经过上述工作，成功研制出两套功能完整、气密性良好的声激励罐体，其装配的结果如图~\ref{fig:16}所示，形成了一套成熟的声激励系统设计、加工和调试流程。解决了密封、电源隔离、电路布局等多个工程技术难题，为后续燃烧实验提供了可靠的声激励源。

\section{温度测量系统}

该部分的目标是建立可重复的二级燃烧室入口温度测量系统，为熄火实验提供可靠的来流温度数据。截至目前，已完成热电偶的多次改造，设计并装配了冷端水冷套，基本解决测温不准和重复性差的问题。

温度采集系统改造与问题排查：针对热电偶测温不准的问题，排查发现其根本原因为热电偶丝过短，冷端受热辐射影响。采取了加长偶丝、安装热挡板、增加陶瓷保护套等措施，初步改善了测温精度。热电偶改造与标定实验通过陶瓷保护套+定位螺母+陶瓷胶的方案，实现了热电偶伸入长度的精确、可重复控制。完成了首次温度标定实验，获得了首批各工况下的入口温度数据。其采集数据结果如图~\ref{fig:32}所示。

\begin{figure}[!htbp]
	\centering
	\includegraphics[width=0.6\textwidth]{image32.png}
	\caption{热电偶记录温度数据}
	\label{fig:32}
\end{figure}

由于进行燃烧实验时，热电偶丝的冷端在横向射流火焰实验台周围，容易受实验台内火焰燃烧时的热辐射影响，因此对其进行改造，增设水冷套。在冷端水冷套设计与装配方面，设计了卡套三通包裹铜管的水冷套结构，用铜管包裹热电偶冷端，外部安装卡套-三通结构通入循环水进行冷却，其结构图如图~\ref{fig:33}所示，下方两端为进水口和出水口，可以对热电偶冷端提供稳定的水冷温度。

\begin{figure}[!htbp]
	\centering
	\includegraphics[width=0.7\textwidth]{image33.jpg}
	\caption{热电偶水冷套装配实物图}
	\label{fig:33}
\end{figure}

该部分建立了一套较为完善的高温温度测量系统，解决了热电偶伸入长度不一致和冷端热辐射干扰等关键问题。设计的冷端水冷套结构简单、成本低廉、冷却效果好，可作为组内高温热电偶测温的标准配置。